\chapter{Fine-Tuning}

Pretrained language models are impressively capable, but they are also
general-purpose: a model trained on internet-scale text can generate
plausible continuations for any prompt, but it cannot follow
instructions, answer questions, or classify sentiment out of the box.
Fine-tuning is the process of adapting a pretrained model to a specific
downstream task or behavior. In this chapter, we cover the core
techniques for fine-tuning: adding classification heads, instruction
fine-tuning with masked loss, attention masks, length-grouped batching,
gradient checkpointing, and parameter-efficient fine-tuning via LoRA.

\section{Classification Head}

Many NLP tasks can be framed as classification: sentiment analysis
(positive/negative), topic categorization, spam detection, and so on.
When adapting a language model for classification, we replace the
language-model head---which projects hidden states to vocabulary-sized
logits---with a smaller linear layer that maps to the number of
classes.

\subsection{From Next-Token Prediction to Classification}

During pretraining, the model learns rich representations of text. The
final hidden states encode semantic information that is useful far
beyond next-token prediction. For classification, we extract a single
vector from the sequence of hidden states and feed it into a
classification head:

\[
  \hat{y} = \text{softmax}\!\left(W_{\text{cls}} \cdot h_{\text{last}} + b_{\text{cls}}\right)
\]

where $h_{\text{last}} \in \mathbb{R}^{d_{\text{model}}}$ is the hidden
state at a chosen position, $W_{\text{cls}} \in \mathbb{R}^{n_{\text{classes}} \times d_{\text{model}}}$,
and $b_{\text{cls}} \in \mathbb{R}^{n_{\text{classes}}}$.

\begin{keyconcept}
  Why the last token? In causal (autoregressive) attention, every token
  can only attend to itself and preceding tokens. The last token has
  attended to \emph{all} preceding tokens, making its hidden state a
  summary of the entire sequence. This is analogous to the
  \texttt{[CLS]} token in BERT, but we don't need a special
  token---the last position naturally accumulates information from the
  full sequence.

  During pretraining, the model is trained so that the last token's
  hidden state predicts the next token given the entire context. This
  forces the last token to encode a compressed representation of
  everything it has seen, which transfers well to classification.
\end{keyconcept}

\begin{pseudocode}[Classification Forward Pass]
def classify(input_ids: Tensor,       # [batch, seq_len]
             model: GPTModel,
             cls_head: Linear,          # Linear(embed_dim, num_classes)
             pad_id: int) -> Tensor:    # [batch, seq_len]
    # Run full model to get hidden states
    hidden = model(input_ids)           # [batch, seq_len, embed_dim]

    # Find the last non-padding token for each sequence
    # attention_mask is 1 for real tokens, 0 for padding
    mask = (input_ids != pad_id)        # [batch, seq_len]
    # argmax finds the last True position (since seq is 0-indexed left-to-right)
    last_idx = mask.int().argmax(dim=1)  # [batch]
    # Gather the hidden state at the last real token
    last_hidden = hidden.gather(           # [batch, embed_dim]
        1, last_idx.unsqueeze(-1).expand(-1, hidden.size(-1))
    ).squeeze(1)
    # Classify
    logits = cls_head(last_hidden)        # [batch, num_classes]
    return logits
\end{pseudocode}

\begin{keyconcept}
  An important detail: when sequences are padded, we must extract the
  hidden state of the \emph{last real token}, not the last position in
  the tensor (which may be a padding token whose hidden state is
  meaningless). The code above uses the attention mask to locate this
  position.
\end{keyconcept}

\subsection{Training the Classification Head}

When fine-tuning for classification, there is a design choice: should
we update the pretrained weights, or freeze them and only train the
classification head?

\begin{itemize}
  \item \textbf{Frozen backbone}: Only $W_{\text{cls}}$ and
        $b_{\text{cls}}$ are updated. Works well when the pretrained
        model already encodes good features and the downstream task is
        simple. Also much faster and uses less memory.
  \item \textbf{Fine-tuning all weights}: The entire model is updated.
        Better for complex tasks that require adapting internal
        representations, but risks overfitting on small datasets and
        requires more memory.
  \item \textbf{Gradual unfreezing}: Start with a frozen backbone and
        progressively unfreeze layers from the top. A common compromise.
\end{itemize}

\section{Instruction Fine-Tuning and Masked Loss}

A pretrained language model generates text unconditionally---it does not
understand the concept of following instructions or engaging in
multi-turn conversation. Instruction fine-tuning teaches the model to
respond to structured prompts by training on examples formatted as
instruction--response pairs.

\subsection{Data Format}

Each training example is formatted as:

\begin{center}
\texttt{<prompt>Instruction here<|separator|>Response here<|eos|>}
\end{center}

The prompt portion provides context and instructions, and the response
is the desired output. During training, we want the model to learn to
\emph{generate the response}, not to copy the prompt. This means we
must compute loss only on the response tokens.

\subsection{Masked Targets}

PyTorch's \texttt{CrossEntropyLoss} supports an \texttt{ignore\_index}
parameter---any position in the target labeled with this index
contributes zero to the loss. By convention, \texttt{ignore\_index =
-100}. We set prompt tokens and padding tokens to $-100$ in the target
tensor:

\[
  \mathcal{L} = -\frac{1}{|\{i : t_i \neq -100\}|}\sum_{i:\, t_i \neq -100} \log P(t_i \mid \text{context})
\]

\begin{pseudocode}[Constructing Masked Targets for Instruction Fine-Tuning]
def build_masked_targets(input_ids: Tensor,       # [batch, seq_len]
                          prompt_lengths: Tensor,   # [batch]
                          pad_id: int,
                          eos_id: int) -> Tensor:   # [batch, seq_len]
    # Start with all positions masked
    targets = torch.full_like(input_ids, -100)  # [batch, seq_len]

    for i in range(input_ids.size(0)):
        prompt_len = prompt_lengths[i].item()
        # Copy response tokens (after prompt separator) into targets
        # Shift by 1: target at position t is the token at position t+1
        targets[i, prompt_len:-1] = input_ids[i, prompt_len + 1:]
        # Mask padding positions in the response
        # Find where real tokens end (excluding eos)
        response = input_ids[i, prompt_len:]
        real_len = (response != pad_id).sum().item()
        # Any remaining positions after real_len are already -100
        # and eos_id position is kept as a valid target

    # Also mask padding positions (where input_ids == pad_id)
    targets[input_ids == pad_id] = -100

    return targets  # [batch, seq_len] with -100 for prompt and padding
\end{pseudocode}

\begin{keyconcept}
  Padding tokens also use \texttt{ignore\_index = -100} so they don't
  contribute to loss. The \texttt{-100} value is PyTorch's convention
  for ``ignore this position.'' This is more memory-efficient than a
  separate binary mask because it reuses the target tensor itself.

  Two types of tokens are masked: (1) prompt tokens, because we don't
  want the model to learn to predict the user's instructions, and (2)
  padding tokens, because they carry no semantic content.
\end{keyconcept}

\section{Attention Masks}

When sequences in a batch have different lengths, shorter sequences are
padded to match the longest sequence in the batch. Attention masks tell
the model which positions contain real tokens and which are padding:

\begin{pseudocode}[Creating and Applying Attention Masks]
def apply_attention_mask(scores: Tensor,       # [batch, num_heads, seq_len, seq_len]
                          input_ids: Tensor,    # [batch, seq_len]
                          pad_id: int) -> Tensor:  # [batch, num_heads, seq_len, seq_len]
    # Create mask: True for real tokens, False for padding
    mask = (input_ids != pad_id)              # [batch, seq_len]

    # Expand mask for multi-head attention:
    # [batch, seq_len] -> [batch, 1, 1, seq_len]
    # Broadcasting will replicate across num_heads and query positions
    mask = mask.unsqueeze(1).unsqueeze(2)      # [batch, 1, 1, seq_len]

    # Apply: set padding positions to -inf before softmax
    # This makes padding tokens get zero attention weight
    scores = scores.masked_fill(~mask, float('-inf'))

    return scores  # [batch, num_heads, seq_len, seq_len]
\end{pseudocode}

\begin{keyconcept}
  The attention mask is expanded from $[\text{batch}, \text{seq}]$ to
  $[\text{batch}, 1, 1, \text{seq}]$ before being applied to the
  $[\text{batch}, \text{heads}, \text{seq}, \text{seq}]$ attention
  scores. Broadcasting replicates it across all heads and all query
  positions.

  During generation (autoregressive decoding), the attention mask must
  \emph{grow} alongside the token sequence---each newly generated token
  gets a mask value of 1 (real). This is handled by concatenating a new
  column of ones to the mask at each generation step.
\end{keyconcept}

\section{Length-Grouped Batching}

Instruction fine-tuning datasets contain sequences of widely varying
lengths. If we naively group random sequences into batches, each batch
must be padded to the length of its longest sequence. A batch containing
one sequence of 900 tokens and seven sequences of 50 tokens wastes
$850 \times 7 = 5{,}950$ padding positions.

\subsection{Grouping by Length}

The solution is to sort sequences by length (with added noise for
randomization) and then group similar-length sequences together into
batches. Short sequences go into one batch, long sequences into
another---each batch is padded only to its own maximum, not the global
maximum.

\begin{pseudocode}[Length-Grouped Batching]
def create_length_grouped_batches(
        examples: list[dict],          # Each has "input_ids" and "target"
        max_tokens: int,                # Max tokens per batch (GPU memory budget)
        noise_factor: float = 0.05      # Noise as fraction of max length
    ) -> list[list[dict]]:
    # Add noise to lengths for randomization within groups
    lengths = [len(ex["input_ids"]) for ex in examples]
    max_len = max(lengths)
    noisy_lengths = [
        l + int(noise_factor * max_len * random())
        for l in lengths
    ]
    # Sort by noisy length
    sorted_examples = sorted(
        zip(noisy_lengths, examples), key=lambda x: x[0]
    )
    examples = [ex for _, ex in sorted_examples]

    # Group into batches respecting token budget
    batches = []
    current_batch = []
    current_max_len = 0

    for ex in examples:
        seq_len = len(ex["input_ids"])
        # Check if adding this example would exceed token budget
        # New batch size * new max length must fit in budget
        new_max_len = max(current_max_len, seq_len)
        new_batch_size = len(current_batch) + 1
        if new_batch_size * new_max_len <= max_tokens:
            current_batch.append(ex)
            current_max_len = new_max_len
        else:
            batches.append(current_batch)
            current_batch = [ex]
            current_max_len = seq_len

    if current_batch:
        batches.append(current_batch)

    # Shuffle batch order to prevent length-biased gradient updates
    shuffle(batches)
    return batches
\end{pseudocode}

\begin{keyconcept}
  Without grouping: a batch containing sequences of length 900 and 50
  results in padding waste (all sequences padded to 900). With grouping,
  short and long sequences go into separate batches, minimizing padding.

  Dynamic batch sizing by token budget (not fixed batch size) is what
  actually prevents OOM: fewer long sequences per batch, more short
  ones. A batch of eight 50-token sequences uses 400 tokens, while a
  batch of two 900-token sequences uses 1{,}800 tokens---both fit within
  a 2{,}048-token budget, but with very different batch sizes.

  The noise factor prevents the model from seeing perfectly sorted data
  every epoch, which could introduce subtle biases.
\end{keyconcept}

\section{Gradient Checkpointing}

During the forward pass, a neural network computes intermediate
activations and stores them in memory for use during backpropagation.
For a Transformer with $L$ layers, these stored activations scale as
$O(L \cdot B \cdot S \cdot D)$, where $B$ is the batch size, $S$ is
the sequence length, and $D$ is the hidden dimension. For large models,
this activation memory can be the dominant memory cost---often larger
than the model weights themselves.

\subsection{How Checkpointing Works}

Gradient checkpointing reduces activation memory by not storing
intermediate activations for every layer. Instead, it stores only
``checkpoint'' activations at selected layer boundaries. During the
backward pass, when gradients are needed for a non-checkpointed layer,
the forward computation from the nearest checkpoint is re-executed to
recompute the required activations.

\begin{pseudocode}[Gradient Checkpointing for Transformer Blocks]
def forward_with_checkpointing(
        x: Tensor,                 # [batch, seq_len, embed_dim]
        blocks: list[TransformerBlock],
        checkpoint_segments: int = 3  # Number of checkpoint boundaries
    ) -> Tensor:
    # Divide blocks into segments; only store activations at boundaries
    blocks_per_segment = len(blocks) // checkpoint_segments

    # Forward pass with checkpointing
    for seg in range(checkpoint_segments):
        start = seg * blocks_per_segment
        end = start + blocks_per_segment
        segment_blocks = blocks[start:end]

        def run_segment(x: Tensor, seg_blocks=segment_blocks) -> Tensor:
            # This function will be re-run during backward if needed
            for block in seg_blocks:
                x = block(x)      # [batch, seq_len, embed_dim]
            return x

        # checkpoint() stores only the input x for this segment
        # During backward, it re-runs run_segment to recompute activations
        x = torch.utils.checkpoint.checkpoint(
            run_segment, x, use_reentrant=False
        )  # [batch, seq_len, embed_dim]

    # Process any remaining blocks (if not evenly divisible)
    for block in blocks[checkpoint_segments * blocks_per_segment:]:
        x = block(x)              # [batch, seq_len, embed_dim]

    return x
\end{pseudocode}

\begin{keyconcept}
  \textbf{Tradeoff}: Gradient checkpointing trades compute for memory.
  By recomputing activations during the backward pass, it saves
  $\sim$55--65\% of activation memory at the cost of $\sim$30\% slower
  training (due to the extra forward passes). On GPT-2 Large (36
  layers):

  \begin{center}
  73 GB $\rightarrow$ 33 GB ($\sim$55\% reduction)
  \end{center}

  The number of checkpoint segments controls the tradeoff: more segments
  means more recomputation but less memory. With a single checkpoint
  per layer (the finest granularity), memory scales as $O(\sqrt{L})$
  instead of $O(L)$. In practice, checkpointing every few layers is
  sufficient.
\end{keyconcept}

\section{LoRA (Low-Rank Adaptation)}

Fine-tuning all parameters of a large language model is expensive: the
full model must be loaded into memory, and optimizer states (Adam
maintains two additional tensors per parameter) can triple the memory
requirement. LoRA (Low-Rank Adaptation) is a parameter-efficient
fine-tuning method that freezes the original weights and adds small
trainable low-rank matrices, dramatically reducing the number of
trainable parameters while maintaining performance comparable to full
fine-tuning.

\subsection{The Low-Rank Decomposition}

LoRA applies a low-rank update to a pretrained weight matrix $W$. The
forward pass becomes:

\[
  h = Wx + \frac{\alpha}{r} BAx
\]

where:
\begin{itemize}
  \item $W \in \mathbb{R}^{d_{\text{out}} \times d_{\text{in}}}$ is the
        frozen pretrained weight matrix
  \item $B \in \mathbb{R}^{d_{\text{out}} \times r}$, initialized to zeros
  \item $A \in \mathbb{R}^{r \times d_{\text{in}}}$, initialized from
        $\mathcal{N}(0, \sigma^2)$ where $\sigma = 0.02$
  \item $r \ll d_{\text{in}}, d_{\text{out}}$ is the LoRA rank
  \item $\alpha$ is the scaling factor; effective scaling is $\alpha / r$
\end{itemize}

\subsection{Why Low Rank Works: A Concrete Example}

Consider a weight matrix $W \in \mathbb{R}^{768 \times 768}$ from
GPT-2 Small, with a LoRA rank of $r = 4$:

\begin{center}
\begin{tabular}{lrr}
\toprule
\textbf{Parameter} & \textbf{Shape} & \textbf{Count} \\
\midrule
$W$ (frozen) & $768 \times 768$ & 589{,}824 \\
$A$ (trainable) & $4 \times 768$ & 3{,}072 \\
$B$ (trainable) & $768 \times 4$ & 3{,}072 \\
LoRA update ($BA$) & $768 \times 768$ (implicit) & --- \\
\bottomrule
\end{tabular}
\end{center}

The update $BA$ produces a $768 \times 768$ matrix, but it is
constrained to lie in a 4-dimensional subspace---it has rank at most 4.
The full matrix $W$ has rank 768 (typically), so the LoRA update can
only modify the output along 4 directions. Yet this is enough: the
pretrained model is already close to a good solution, and fine-tuning
requires only small adjustments that tend to be low-rank.

The parameter savings are stark: 6{,}144 trainable parameters for the
LoRA adapter versus 589{,}824 for full fine-tuning of this single
matrix---a $\sim$96$\times$ reduction. Across all attention weight
matrices in GPT-2 Small (12 layers $\times$ 4 matrices each), the
total LoRA parameters are roughly 294K versus 780M for full
fine-tuning.

\subsection{Forward Pass}

\begin{pseudocode}[LoRA Linear Layer Forward Pass]
class LoRALinear(nn.Module):
    def __init__(self, in_dim: int, out_dim: int, rank: int, alpha: float):
        super().__init__()
        # Frozen pretrained weight
        self.W = nn.Linear(in_dim, out_dim, bias=False)
        self.W.weight.requires_grad_(False)

        # LoRA low-rank matrices
        self.lora_A = nn.Parameter(torch.randn(rank, in_dim) * 0.02)
        self.lora_B = nn.Parameter(torch.zeros(out_dim, rank))
        self.scaling = alpha / rank

    def forward(self, x: Tensor) -> Tensor:  # x: [batch, seq_len, in_dim]
        # Pretrained path (frozen)
        h_pretrained = self.W(x)              # [batch, seq_len, out_dim]

        # LoRA path: x -> A -> B (low-rank bottleneck)
        # Process sequence dimension with matrix operations
        h_lora = x @ self.lora_A.T             # [batch, seq_len, rank]
        h_lora = h_lora @ self.lora_B.T         # [batch, seq_len, out_dim]

        # Combine with scaling
        return h_pretrained + self.scaling * h_lora  # [batch, seq_len, out_dim]
\end{pseudocode}

\begin{keyconcept}
  \textbf{Why $B$ is zero-initialized}: At the start of fine-tuning,
  $BA = 0$ so $h = Wx$, meaning the model behaves exactly like the
  pretrained model. The LoRA update grows from zero, ensuring a smooth
  transition from the pretrained model's behavior.

  If $B$ were randomly initialized, the initial output would be $h =
  Wx + \frac{\alpha}{r}BAx$, introducing a large random perturbation
  that could destabilize training and temporarily degrade the model's
  capabilities.

  \textbf{Practical impact}: LoRA reduces trainable parameters from
  $\sim$780M to $\sim$6M on GPT-2 Large. The main savings come from
  eliminating optimizer states (Adam stores 2 extra tensors per
  parameter) and gradients for frozen weights. With LoRA, optimizer
  state memory scales with the number of \emph{trainable} parameters,
  not the total model size.

  \textbf{Merging for deployment}: At inference time, the LoRA weights
  can be merged into the original weights: $W' = W + \frac{\alpha}{r}
  BA$. After merging, there is no additional cost compared to a fully
  fine-tuned model---same latency, same memory.
\end{keyconcept}

\subsection{Which Layers to Apply LoRA To}

LoRA is typically applied to the attention weight matrices ($W_Q$,
$W_K$, $W_V$, $W_O$) and sometimes the feed-forward layers. Empirical
results show that:

\begin{itemize}
  \item Applying LoRA to all four attention matrices ($W_Q, W_K, W_V,
        W_O$) generally outperforms applying it to only $W_Q$ and
        $W_V$.
  \item Applying LoRA to feed-forward layers in addition to attention
        brings marginal improvement at the cost of more parameters.
  \item The rank $r$ is typically set between 4 and 64. Ranks of 8--16
        often provide a good balance between expressiveness and
        efficiency for instruction fine-tuning.
\end{itemize}

\section*{Review Questions}

\begin{reviewquestion}[1]
  Why does instruction fine-tuning use \texttt{ignore\_index = -100}
  instead of a separate mask tensor? What are the two types of tokens
  that should be masked out?
\end{reviewquestion}

\begin{reviewquestion}[2]
  Explain how length-grouped batching reduces padding waste. Why do
  we add noise to sequence lengths before sorting?
\end{reviewquestion}

\begin{reviewquestion}[3]
  LoRA initializes $B$ to zeros. Explain why this is important for
  fine-tuning. What would happen if $B$ were randomly initialized
  instead?
\end{reviewquestion}

\begin{reviewquestion}[4]
  In classification fine-tuning, why do we extract the hidden state of
  the last real token rather than the last position in the tensor? What
  would happen if we used the last position in a padded sequence?
\end{reviewquestion}

\begin{reviewquestion}[5]
  Explain the tradeoff made by gradient checkpointing. If you
  checkpoint every $k$ layers in a 24-layer model, approximately how
  many times is each layer's forward pass executed during one training
  iteration?
\end{reviewquestion}

\begin{reviewquestion}[6]
  For a $1024 \times 768$ weight matrix with LoRA rank $r = 8$ and
  $\alpha = 16$, compute the number of trainable parameters in the LoRA
  adapter. What is the effective scaling factor $\alpha / r$?
\end{reviewquestion}

\begin{reviewquestion}[7]
  Why does causal attention make the last token's hidden state a good
  representation for classification? Contrast this with BERT-style
  bidirectional attention.
\end{reviewquestion}

\begin{reviewquestion}[8]
  When building masked targets for instruction fine-tuning, the targets
  are shifted by one position relative to the input. Explain why this
  shift is necessary.
\end{reviewquestion}

\begin{reviewquestion}[9]
  Describe how dynamic batch sizing based on a token budget prevents
  out-of-memory errors during fine-tuning. Why can't we just use a
  fixed batch size?
\end{reviewquestion}

\begin{reviewquestion}[10]
  After LoRA fine-tuning, we can merge the LoRA weights into the
  original weights via $W' = W + \frac{\alpha}{r}BA$. What is the
  advantage of merging at deployment time? Is there any difference in
  inference latency between merged and unmerged LoRA?
\end{reviewquestion}

\begin{reviewquestion}[11]
  Gradient checkpointing reduces memory from $O(L)$ to $O(\sqrt{L})$
  for $L$ layers (with per-layer checkpointing). Derive why the memory
  scales as $O(\sqrt{L})$ when every layer is a checkpoint boundary.
\end{reviewquestion}

\begin{reviewquestion}[12]
  In the LoRA forward pass, the computation goes through a bottleneck:
  $x \to A \to B \to$ output, where $A$ is $r \times d_{\text{in}}$
  and $B$ is $d_{\text{out}} \times r$. Explain why the rank $r$
  controls the expressiveness of the LoRA update. What happens as $r
  \to d_{\text{in}}$?
\end{reviewquestion}

\section*{Problems}

\begin{problem}[1]
  \textbf{Implement a classification fine-tuning loop.} Given a
  pretrained GPT model, a classification head
  (\texttt{nn.Linear(embed\_dim, num\_classes)}), and a dataset of
  (text, label) pairs, implement the training loop that:
  (a) extracts the last real token's hidden state, (b) computes
  classification loss, (c) optionally freezes the backbone. Compare
  validation accuracy with a frozen backbone vs.\ fine-tuning all
  weights on a binary sentiment classification task.
\end{problem}

\begin{problem}[2]
  \textbf{Implement masked targets for instruction fine-tuning.}
  Write a function that takes a batch of tokenized instruction--response
  pairs (with a separator token and padding) and returns the target
  tensor with prompt and padding positions set to $-100$. Verify your
  implementation by checking that: (a) the loss on a batch equals the
  loss computed only on response tokens, and (b) padding tokens
  contribute zero to the loss.
\end{problem}

\begin{problem}[3]
  \textbf{Implement LoRA adapter layers.} Create a
  \texttt{LoRALinear} module that wraps an existing \texttt{nn.Linear}
  layer and adds the low-rank $A$ and $B$ matrices with configurable
  rank $r$ and scaling $\alpha$. Replace the attention weight matrices
  in a GPT model with LoRA-wrapped versions. Verify that at
  initialization, the LoRA model produces identical outputs to the
  original model (since $B$ is zero-initialized).
\end{problem}

\begin{problem}[4]
  \textbf{Implement length-grouped batching.} Given a list of
  tokenized sequences of varying lengths, implement a dataloader that
  sorts sequences by length (with a noise factor), groups them into
  batches respecting a maximum token budget, and pads each batch to its
  own maximum length. Measure the average padding waste per batch
  compared to random batching.
\end{problem}

\begin{problem}[5]
  \textbf{Implement gradient checkpointing.} Wrap each
  \texttt{TransformerBlock} in a GPT model with
  \texttt{torch.utils.checkpoint.checkpoint()} and measure: (a) peak
  GPU memory with and without checkpointing, (b) training throughput
  (tokens per second) with and without checkpointing. Report the
  memory savings and the throughput reduction.
\end{problem}